\section{Requirements for an MCP “middle layer” in public-sector settings (derived from the case)}

These requirements are written to be \textbf{testable}, \textbf{auditable}, and aligned to public-sector realities (mixed sensitivity, mixed data quality, scrutiny, and real-time questioning).

\textbf{Primary evidence used:} transcript + tool trace + trace-derived forensic analyses \cite{artefact_claude_stopped_after_ngd_features_2026_02_18,artefact_codex_trace_jsonl_2026_02_18,artefact_floor_question_analysis_2026_02_19,artefact_forensic_study_mcp_geo_2026_02_19}.

\textbf{Standards and guidance anchors} used for requirements framing include MCP spec, provenance/catalogue standards, and UK assurance guidance. \cite{mcp_spec_2025_11_25,prov_o,prov_dm,dcat3,ncsc_secure_ai_guidelines,ico_ai_data_protection,data_ai_ethics_framework,five_safes_framework,algorithmic_transparency_recording_standard,owasp_llm_top10}

Each requirement includes an \textbf{acceptance test}.

---

\subsection*{\thesection.R01 Intent routing: map questions to an explicit workflow}
\textbf{Requirement:} The server MUST classify each request into an intent class (e.g., \textit{named-area lookup}, \textit{boundary resolution}, \textit{indicator scan}, \textit{detailed geometry extract}) and route to an explicit workflow with safe defaults. \cite{mcp_spec_2025_11_25,owasp_llm_top10,ncsc_secure_ai_guidelines}

\textbf{Acceptance tests:}
\begin{itemize}[leftmargin=*]
\item Given “peatland site survey … Bowland”, the server selects \textit{environmental survey} intent and triggers \textit{boundary resolution} before any heavy NGD geometry calls.
\item The chosen workflow is recorded in the audit log.
\end{itemize}

\subsection*{\thesection.R02 Boundary catalogue: resolve named landscapes beyond administrative geographies}
\textbf{Requirement:} The server MUST support authoritative boundaries for common designated areas (at minimum: National Landscapes/AONB, National Parks, SSSI; configurable by nation) and MUST label when a proxy boundary is used. \cite{dcat3,data_ai_ethics_framework,algorithmic_transparency_recording_standard}

\textbf{Acceptance tests:}
\begin{itemize}[leftmargin=*]
\item “Forest of Bowland” resolves to an authoritative National Landscape polygon when available.
\item If not available, the output includes \texttt{boundaryType=proxy}, \texttt{proxyReason}, and the exact proxy dataset used.
\end{itemize}

\subsection*{\thesection.R03 “Counts-first” default mode for large-area survey questions}
\textbf{Requirement:} For survey-style questions, the server MUST default to a low-payload indicator scan: \cite{owasp_llm_top10,ncsc_secure_ai_guidelines}
\begin{itemize}[leftmargin=*]
\item \texttt{resultType=hits},
\item \texttt{includeGeometry=false},
\item minimal fields,
\item and spatial tiling if the extent is large.
\end{itemize}

\textbf{Acceptance tests:}
\begin{itemize}[leftmargin=*]
\item A new “survey” request returns an indicator summary without any tool output exceeding a configured size threshold (e.g., 50 kB).
\item Geometry is only fetched after the user or workflow explicitly requests it.
\end{itemize}

\subsection*{\thesection.R04 Payload guardrails with progressive disclosure}
\textbf{Requirement:} The server MUST enforce payload budgets and progressive disclosure: \cite{owasp_llm_top10,ncsc_secure_ai_guidelines}
\begin{itemize}[leftmargin=*]
\item clamp \texttt{limit} to upstream maxima,
\item warn on large bboxes,
\item split requests into tiles/pages,
\item and never stream multi-hundred-kilobyte feature collections inline by default.
\end{itemize}

\textbf{Acceptance tests:}
\begin{itemize}[leftmargin=*]
\item Requests with \texttt{limit>100} are clamped and a warning is returned.
\item Any response >200 kB is converted into a summary + handle (see R5), not inline JSON.
\end{itemize}

\subsection*{\thesection.R05 Resource-backed result delivery for large datasets}
\textbf{Requirement:} When tool outputs exceed a threshold, the server MUST store full results server-side (or in an approved blob store) and return: \cite{mcp_spec_2025_11_25,owasp_llm_top10}
\begin{itemize}[leftmargin=*]
\item a compact summary (counts, bbox, sample IDs), and
\item a stable \texttt{resource://} handle for paging/export.
\end{itemize}

\textbf{Acceptance tests:}
\begin{itemize}[leftmargin=*]
\item The “large geometry” scenario returns a handle and a summary; follow-up “fetch next page” uses the handle without recomputing earlier pages.
\item Resources include TTL and access controls appropriate to data sensitivity.
\end{itemize}

\subsection*{\thesection.R06 Correct and explicit counting semantics}
\textbf{Requirement:} The server MUST provide unambiguous counters: \cite{algorithmic_transparency_recording_standard,prov_dm}
\begin{itemize}[leftmargin=*]
\item \texttt{numberReturned} MUST equal the number of features in \texttt{features[]},
\item \texttt{count} MUST be explicitly defined (returned count vs total matched),
\item and when local filtering is used, the response MUST state \texttt{filterApplied=local} and the scan scope/budget.
\end{itemize}

\textbf{Acceptance tests:}
\begin{itemize}[leftmargin=*]
\item Empty result sets always return \texttt{numberReturned=0}.
\item A regression test covers \texttt{filter}+\texttt{resultType=results} and confirms counters remain consistent.
\end{itemize}

\subsection*{\thesection.R07 Explicit filter model: upstream CQL vs local post-filter}
\textbf{Requirement:} The tool schema MUST distinguish: \cite{prov_dm,algorithmic_transparency_recording_standard,owasp_llm_top10}
\begin{itemize}[leftmargin=*]
\item upstream filters (e.g., CQL, filter-lang),
\item and local post-filters.
\end{itemize}
If local post-filtering is used, the server MUST prevent global claims unless all pages are scanned (or it explicitly marks the result as partial).

\textbf{Acceptance tests:}
\begin{itemize}[leftmargin=*]
\item For a local filter with paging, the response contains \texttt{partialScan=true} unless all pages are scanned within the budget.
\item The report generator refuses to state “no peat present” unless the scan completeness is 100\% or an authoritative peat dataset is used.
\end{itemize}

\subsection*{\thesection.R08 Strict geometry input validation (no silent downgrades)}
\textbf{Requirement:} Any supplied polygon/geometry constraint MUST be validated and either applied or rejected. The server MUST NOT silently ignore a geometry constraint. \cite{ncsc_secure_ai_guidelines,owasp_llm_top10}

\textbf{Acceptance tests:}
\begin{itemize}[leftmargin=*]
\item If \texttt{polygon} is provided as a string (wrong type), the tool returns \texttt{INVALID\_INPUT} with a helpful message.
\item If both bbox and polygon are supplied, the server reports which constraint is applied and why.
\end{itemize}

\subsection*{\thesection.R09 Resilience: adaptive retries and degrade plan for upstream errors}
\textbf{Requirement:} On upstream timeouts or transient errors, the server MUST: \cite{ncsc_secure_ai_guidelines,owasp_llm_top10}
\begin{itemize}[leftmargin=*]
\item retry with exponential backoff (idempotent calls only),
\item and apply a “degrade plan” (reduce bbox, drop geometry, lower limit) before failing.
\end{itemize}

\textbf{Acceptance tests:}
\begin{itemize}[leftmargin=*]
\item In a simulated timeout, the second attempt uses reduced cost parameters and succeeds, with an audit entry describing the degrade plan.
\item The final user output includes which calls were retried and which degradations were applied.
\end{itemize}

\subsection*{\thesection.R10 Provenance-by-default output blocks}
\textbf{Requirement:} Every user-facing answer MUST include a provenance block: \cite{prov_o,prov_dm,algorithmic_transparency_recording_standard,dcat3}
\begin{itemize}[leftmargin=*]
\item datasets used (name, version/date, licensing),
\item boundary source + type (authoritative/proxy),
\item query parameters (bbox/polygon, paging completeness),
\item and explicit limitations.
\end{itemize}

\textbf{Acceptance tests:}
\begin{itemize}[leftmargin=*]
\item Generated answers always include a provenance block, even in “quick answer” mode.
\item Removing the provenance block fails CI for the report generator.
\end{itemize}

\subsection*{\thesection.R11 Sensitivity guardrails and policy-aware redaction}
\textbf{Requirement:} The server MUST support data-classification-aware behaviour: \cite{ico_ai_data_protection,five_safes_framework,data_ai_ethics_framework,ncsc_secure_ai_guidelines}
\begin{itemize}[leftmargin=*]
\item prevent accidental retrieval/exposure of sensitive attributes,
\item redact or aggregate by default,
\item and require explicit privilege and justification for sensitive layers.
\end{itemize}

\textbf{Acceptance tests:}
\begin{itemize}[leftmargin=*]
\item A restricted dataset cannot be accessed without a declared clearance token/role.
\item Responses from restricted datasets are automatically summarised/aggregated and audited.
\end{itemize}

\subsection*{\thesection.R12 Audit-ready observability (trace → report)}
\textbf{Requirement:} The server MUST emit structured audit logs that can be deterministically replayed into a timeline report: \cite{algorithmic_transparency_recording_standard,prov_o,mcp_spec_2025_11_25,ncsc_secure_ai_guidelines}
\begin{itemize}[leftmargin=*]
\item tool calls (args, status, latency, sizes),
\item routing decisions,
\item warnings/guardrails triggered,
\item and resource handles created.
\end{itemize}

\textbf{Acceptance tests:}
\begin{itemize}[leftmargin=*]
\item A report generator can rebuild a timeline table identical (within tolerances) to the runtime record.
\item Correlation IDs link host messages to tool calls, reducing host–server mismatch risk.
\end{itemize}

\textit{Provenance reference:} W3C PROV-O. \cite{prov_o}

\textit{Catalogue reference:} W3C DCAT v3. \cite{dcat3}
